% =====================================================================
%  Solving Coalgebraic Equations in the Pure $\lambda$-Calculus
% ---------------------------------------------------------------------
%  Shared body. This file carries NO \documentclass: it is \input by two
%  entry points that supply the class (mimicking ../inheritance-calculus):
%    * preprint.tex   -- arXiv / preprint  (acmart, nonacm)
%    * submission.tex -- POPL submission   (acmart, review, anonymous)
%  Build both with: latexmk  (see latexmkrc)
% ---------------------------------------------------------------------
%  STATUS: working draft.
%  * Placeholders are marked \TODO{...} and render in red.
%  * Every bibliography entry is UNVERIFIED: author/title are best-effort
%    recollections; years, venues, volumes, and pages are marked [verify].
%    Do not cite from this file without checking each entry.
%  * Proofs are sketches; load-bearing gaps are flagged with \TODO.
% =====================================================================

%% acmart already loads amsmath, amssymb, amsthm, xcolor, and hyperref;
%% only mathtools needs to be added here.
\usepackage{mathtools}

%% listings renders the machine-generated Python in the compiler appendix.
\usepackage{listings}
\lstset{
  basicstyle=\ttfamily\small,
  breaklines=true,
  breakatwhitespace=false,
  columns=fullflexible,
  keepspaces=true,
  showstringspaces=false,
  language=Python,
  xleftmargin=1em,
  aboveskip=0.4\baselineskip,
  belowskip=0.4\baselineskip,
}

% --- draft placeholder macro ------------------------------------------
\newcommand{\TODO}[1]{\textcolor{red}{[\textbf{TODO:} #1]}}
\newcommand{\citeTODO}[1]{\textcolor{red}{[\textbf{cite:} #1]}}

% --- general notation --------------------------------------------------
\newcommand{\Set}{\mathbf{Set}}
\newcommand{\Beh}{B}                         % behaviour functor B = M F
\newcommand{\rfp}{\varrho}                   % rational fixpoint
\newcommand{\behav}{\mathsf{beh}}            % behaviour map into the final coalgebra
\newcommand{\outm}{\mathsf{out}}             % structure map: the iteration of the equation morphism
\newcommand{\stepm}{e}                       % equation morphism (the map e)
\newcommand{\Delaym}{\mathsf{Delay}}
\newcommand{\Pow}{\mathcal{P}}
\newcommand{\iter}[1]{#1^{\dagger}}          % Elgot iteration

% --- lambda-instance notation (used only in the HNF bridge/examples) ---
\newcommand{\Sh}{\mathsf{Sh}}
\newcommand{\Lamsh}{\mathsf{Lam}}
\newcommand{\Appsh}{\mathsf{App}}
\newcommand{\Varsh}{\mathsf{Var}}
\newcommand{\dlam}{\mathord{\downarrow}\lambda}
\newcommand{\dhead}{\mathord{\downarrow}\mathsf{h}}
\newcommand{\darg}{\mathord{\downarrow}\mathsf{a}}
\newcommand{\subt}[2]{#1\!\downarrow\!#2}
\newcommand{\reduces}{\twoheadrightarrow}
\newcommand{\bred}{\rightarrow_{\beta}}
\newcommand{\present}{\mathtt{present}}

% --- theorem environments ---------------------------------------------
%% acmart provides theorem/lemma/proposition/corollary (acmplain) and
%% definition/example (acmdefinition), all sharing the `theorem' counter.
%% Only `remark' is missing, so we add it on the same shared counter.
\AtBeginDocument{%
  \theoremstyle{acmdefinition}%
  \newtheorem{remark}[theorem]{Remark}%
}

\begin{document}

\title{Solving Coalgebraic Equations in the Pure $\lambda$-Calculus}

\author{Bo Yang}
\affiliation{%
  \institution{Figure AI Inc.}
  \city{San Jose}
  \state{California}
  \country{USA}
}
\email{yang-bo@yang-bo.com}
\thanks{The author completed this work independently and outside the scope of employment at Figure AI. No Figure AI time, funding, or other resources were used in support of this work.}

\begin{abstract}
  Representing and transforming cyclic and infinite data in the pure $\lambda$-calculus normally requires an added recursion construct, a \texttt{letrec}, a $\mu$-binder, or a built-in $Y$ for graph reduction: such data has no finite normal form, so ordinary reduction unfolds it forever instead of folding it into a finite cyclic graph, and an ordinary recursive function over it does the same. We show no extension is needed: a term keeps the meaning its reduction defines, which for weak-head reduction is its standard lazy (L\'evy--Longo) semantics, and which we represent as a finite cyclic graph when it is rational. The insight is that such data is the unique solution of a coalgebraic equation, and that any coalgebraic equation can be solved directly, by tabling: when a state recurs, the solver folds it into a back edge, so rational data, that with finitely many distinct subtrees, comes out as a finite cyclic graph. The solver halts exactly on the inputs with finitely many reachable states, each reached by a terminating reduction, a sub-fragment of the rational data, since not every rational datum reaches finitely many states; the rational fragment itself is the ceiling for finite cyclic graphs, as data beyond it has no such graph, so no solution-preserving procedure returns one on more. The algorithm's first application is the pure $\lambda$-calculus itself: head normalisation is such an equation, and because its states are first-order graph nodes rather than functions their equality is decidable where $\beta$-equality is not, so the solver is an interpreter. It decides $\Omega$ as $\bot$ in finite time and evaluates $\mathtt{map}\,\mathtt{succ}\,(Y\,(\mathtt{cons}\;0))$, an ordinary map over the cyclic stream of zeros, to the cyclic stream of ones as a finite cyclic graph, with no construct added. Because any computable coalgebra Scott-encodes as a $\lambda$-term, this one interpreter solves them all on their rational fragments, making the pure $\lambda$-calculus a domain-specific language whose memoization and transposition tables fall out of node identity: Datalog least Herbrand models, graph reachability and Andersen points-to analysis, minimax search, and dynamic programming. A second example coalgebra is compilation. Its equation sends each sub-term of a source program, by specialization rules, to a node of a Python AST over that sub-term's own sub-terms, and its solution is the program optimised to that AST. A compiler computes it by the same tabling, compiling each shared sub-term once into a finite graph, and is itself a $\lambda$-term the interpreter solves, so program analysis is expressed as evaluation, and it compiles its own source in a self-hosting bootstrap.
\end{abstract}

\keywords{coalgebra, coalgebraic equations, rational fixpoint, algebraic compactness,
  structure as identity, least fixpoint, rational trees, tabling, head reduction, Datalog,
  program analysis, domain-specific language,
  \texorpdfstring{$\lambda$-calculus}{lambda-calculus},
\texorpdfstring{B\"ohm trees}{Bohm trees}}

\maketitle

\section{Introduction}\label{sec:intro}

Ordinary reduction in the pure $\lambda$-calculus never builds a cyclic data structure. Consider
the stream of zeros, the pure term $r = Y\,(\mathtt{cons}\,0)$, with the fixed-point combinator
$Y$ and $\mathtt{cons}\,h\,t$ the stream $t$ with $h$ prepended. The term $r$
reduces to $\mathtt{cons}\,0\,r$, so its tail is $r$ again: the stream that loops on itself. Yet
$r$ has no finite normal form: ordinary reduction unfolds it to $0,0,0,\dots$ forever and never
folds it into a finite circular graph. To obtain a finite cycle one reaches outside the calculus
for an added recursion construct, a \texttt{letrec}, say, that names the
loop~\cite{ariola1997-explicit-recursion}.

Creating the cycle is only half of the difficulty. Take the textbook $\mathtt{map}$, which
applies a function to every element of a list and is itself an ordinary pure $\lambda$-term
(Appendix~\ref{app:dsl}). On the cyclic stream $r$, the term $\mathtt{map}\;\mathtt{succ}\;r$
(with $\mathtt{succ}$ adding one) reduces to $\mathtt{cons}\,1\,(\mathtt{map}\;\mathtt{succ}\;r)$,
whose inner $\mathtt{map}\;\mathtt{succ}\;r$ is the same term again, so an ordinary evaluator
walks the cycle forever and never folds the result back into a cycle. Nothing in the term records
that $r$ was already visited. To make the
traversal terminate the function
must itself become cycle-aware: it has to remember the nodes already visited, which needs a
mutable set or
dictionary, and it has to recognise a revisited node, which needs a reference-equality primitive
that compares nodes by pointer identity. That primitive is what a pure language forbids:
reference equality tells apart structures of equal value, so it breaks referential transparency,
the rule that equal values may be substituted freely.

The standard account therefore pays twice, an extension to create the cycle and an impurity to
transform it.

The obstruction is not special to infinite data. Edit distance, the fewest single-character
insertions, deletions, and substitutions that turn one string into another, has the textbook
recurrence~\cite{wagner1974-string-correction} on the strings' suffixes: reading head by head, equal heads recurse on both tails at
no cost, unequal heads pay one and take the best of a substitution, a deletion, and an insertion,
and once a string is empty the cost is the length of what remains of the other. Written as a pure
$\lambda$-term, with each string a list that applies a head/tail handler $\lambda h.\lambda t$ or,
when empty, a default,\footnote{Data here are encoded as their own case analysis (the Scott and
Church encodings~\cite{mogensen1992-self-interpretation}): a list applied to two arguments runs the first on its head and tail or returns
the second when empty, and a boolean applied to two arguments returns the first if true and the
second if false, so $\mathtt{eq}\,h_a\,h_b$ applied to the two terms that follow it selects one of
them.} and with $\mathtt{eq}$ comparing two characters, $\mathtt{min}_3$ the three-way minimum,
$\mathtt{succ}$ adding one, and $\mathtt{len}$ the length:
\[
  \begin{aligned}
    &\mathtt{ed} = Y\,\big(\lambda\mathit{ed}.\,\lambda a.\,\lambda b.\\
    &\quad a\,\big(\lambda h_a.\lambda t_a.\\
    &\qquad b\,\big(\lambda h_b.\lambda t_b.\\
    &\qquad\quad (\mathtt{eq}\,h_a\,h_b)
        &&\text{compare the two heads}\\
    &\qquad\qquad (\mathit{ed}\,t_a\,t_b)
        &&\text{equal: recurse on both tails, no cost}\\
    &\qquad\qquad \mathtt{succ}\,\big(\mathtt{min}_3\,
        \underbrace{(\mathit{ed}\,t_a\,t_b)}_{\text{substitute}}\,
        \underbrace{(\mathit{ed}\,t_a\,b)}_{\text{delete}}\,
        \underbrace{(\mathit{ed}\,a\,t_b)}_{\text{insert}}\big)\big)\\
    &\qquad (\mathtt{len}\,a)\big)
        &&\text{$b$ empty: delete all of $a$}\\
    &\quad (\mathtt{len}\,b)\big).
        &&\text{$a$ empty: insert all of $b$}
  \end{aligned}
\]
This term computes the right distance, but its naive evaluation is exponential, the three
recursive calls re-expanding overlapping suffix pairs. The subproblems are pairs of suffixes
$(t_a,t_b)$, and a suffix is a shared sub-node of the inputs, not a copy, so if each repeated
pair were computed once the work would be the $(m{+}1)(n{+}1)$ distinct pairs for inputs of
lengths $m$ and $n$, the classic $O(mn)$ dynamic-programming table. But ordinary reduction has
no way to notice that a pair has recurred, so it recomputes. Like the cyclic $\mathtt{map}$, the obvious pure program is correct yet
unusable: one diverges, the other explodes.

The resolution adds nothing to the calculus. We give an interpreter, and a compiler, for the
pure $\lambda$-calculus under which these very programs run efficiently, their meaning unchanged:
$Y\,(\mathtt{cons}\,0)$ evaluates to a finite cyclic graph, $\mathtt{map}\,\mathtt{succ}$ over it
to the finite circle of ones, and $\mathtt{ed}$ in $O(mn)$ time and space. An implementation
accompanies the paper\anon[ as supplementary material]{, openly
available at \url{https://github.com/Atry/MIXINv2}~\cite{colambda}}. The one idea is to read a
program as a \emph{coalgebraic equation}, a state exposing one
layer of data over its successor states, and let the $\lambda$-calculus be its \emph{solver}: it
tables the states it has seen and folds a recurrence on the decidable identity of first-order
graph nodes, the syntax a state is built from, where the only identity on the functions they
compute, $\beta$-equality, is undecidable. The recognition is the solver's, not an operation a
term performs, so the calculus stays pure. And the solution is not checked against the meaning,
it \emph{is} the meaning: what a program denotes is by definition the tree its equation unfolds
to, so computing the solution cannot change the semantics (Section~\ref{sec:bridge}).

The construction is not special to streams, strings, or the $\lambda$-calculus. The tabling
needs only states that expose layers and name successor states (Section~\ref{sec:coalgebras});
it never asks how a state produces its layer. Any reduction function supplies such states: a
rule that at each state either exposes data or takes an internal step, the layer being whatever
the internal steps eventually expose, with the $\lambda$-calculus's own reduction, stopped as
soon as the outermost shape of the term is settled, as one instance
(Section~\ref{sec:reductions}). So for any reduction
function, wherever the solver halts it returns exactly the semantics that function defines. It
halts on the data that reaches finitely many states, a strict sub-fragment of the rational data,
those with finitely many distinct sub-parts; and beyond the rational data no solution-preserving
procedure halts with a finite cyclic graph at all, since there is none to return
(Theorem~\ref{thm:optimal}). And since data such as lists and trees are themselves pure
$\lambda$-terms (the Scott encoding), every such computable system is again a $\lambda$-term, and
one interpreter solves them all on their rational fragments, turning the pure $\lambda$-calculus into a small domain-specific
language whose memoization falls out of node identity: Datalog queries, graph reachability and
points-to analysis, game-tree (minimax) search, and dynamic programming
(Section~\ref{sec:application}, Appendix~\ref{app:dsl}).

\paragraph{Contributions}
\begin{itemize}
  \item \emph{The construction.} A solver for any coalgebra by \emph{tabling}: follow its
    one-step map a layer at a time and reuse a tabled state whenever one recurs, so the back
    edges that fold a cycle cost nothing extra (Sections~\ref{sec:coalgebras},
    \ref{sec:tabling}). Any reduction function presents such a coalgebra through the least
    fixpoint of its internal steps, under which an unproductive loop, one returning to its
    start without exposing a layer, is decided as $\bot$ (undefined) in finite time
    (Section~\ref{sec:reductions}; Definition~\ref{def:out}). The construction is parametric in
    an effect (the internal process, for instance $\beta$-reduction) and a signature (the shape
    of each layer).
  \item \emph{The theory.} We prove the solver \emph{sound} (the finite graph it returns
    unfolds to exactly the solution, Theorem~\ref{thm:correctness}), \emph{bounded at the rationality ceiling} (no
    solution-preserving procedure clears the rational fragment, and the solver halts on a decidable
    sub-fragment of it, the data reaching finitely many states, Theorems~\ref{thm:termination},
    \ref{thm:optimal}), and the equation
    \emph{uniquely solved} (least equals greatest, so induction and coinduction coincide and
    the least fixpoint does real work only in deciding an unproductive loop as $\bot$,
    Theorem~\ref{thm:unique}, Remarks~\ref{rem:compact}, \ref{rem:leastessential}).
  \item \emph{Instances.} Head normalisation of the pure $\lambda$-calculus is one such equation,
    the interpreter (Sections~\ref{sec:bridge}, \ref{sec:application}), on which an ordinary
    $\mathtt{map}$ folds a circular list into a circular list and a naive walk terminates where
    evaluation diverges; because any computable coalgebra Scott-encodes, the one interpreter is also a small
    domain-specific language solving Datalog, graph reachability and points-to analysis, minimax
    search, and dynamic programming (Appendix~\ref{app:dsl}); and compilation is itself such an
    equation, the same interpreter solving it into a self-hosting compiler to Python
    (Section~\ref{sec:application}).
\end{itemize}

% =====================================================================
\section{The Method}\label{sec:method}
% =====================================================================
The method has three independent parts, presented general to specific. A
\emph{coalgebra}~\cite{rutten2000-universal-coalgebra} is
the object we solve, and its solution is what it means (Section~\ref{sec:coalgebras}). The
\emph{approximation order} makes the solution well defined, as the unique fixpoint of a guarded
functional (Section~\ref{sec:order}). \emph{Tabling} computes it (Section~\ref{sec:tabling}).
None of the three mentions how the coalgebra arose; a \emph{reduction function} is one way to
present a coalgebra, through the single least fixpoint of the paper, and the
$\lambda$-calculus enters only there, as an instance (Section~\ref{sec:reductions}).

\subsection{Coalgebras and Their Solutions}\label{sec:coalgebras}

\begin{definition}[Signature and effect]\label{def:signature}
  A \emph{signature} $F$ is a set of operations, each with a finite arity, with decidable and
  hashable equality of operations; $F X$ is the set of layers $\sigma(x_1,\dots,x_k)$ with
  $\sigma$ an operation of arity $k$ and $x_i \in X$. An \emph{effect} $M$ assigns to each set
  $A$ a set $M A$ of computations over $A$ together with a divergent computation
  $\bot_{MA}$; the deterministic effect is $\Delaym$, a computation being either a value or
  divergence, and it is the only effect the examples of this paper use, so a reader may take
  $M(F X) = F X \cup \{\bot\}$ throughout. Richer effects (non-deterministic, probabilistic,
  stateful) fit the same architecture, the silent iteration becoming that of a complete Elgot
  monad \cite{adamek2006-iterative-algebras, goncharov2016-complete-elgot-monads}; this paper
  develops and proves only the deterministic case, on which the metric argument of
  Theorem~\ref{thm:unique} relies (a layer exposes a single operation at the root).
\end{definition}

\begin{definition}[Coalgebra]\label{def:coalgebra}
  A \emph{coalgebra} is a set $X$ of \emph{states}, with decidable and hashable identity,
  together with a map
  \[
    \outm \colon X \longrightarrow M(F X).
  \]
  At each state, $\outm$ either exposes one layer of data, an operation applied to successor
  states, or is the divergent computation $\bot$.
\end{definition}

A coalgebra poses an equation. Each state names a value: one layer of data over the values of
its successor states, which are again states. To \emph{solve} the equation is to compute, for
each state, the tree this naming unfolds to: expose a layer, then recursively the trees of the
successors, with a $\bot$ leaf where $\outm$ diverges. We write $\behav(x)$ for the tree at
state $x$ and call $\behav$ the \emph{solution} of the coalgebra; Section~\ref{sec:order} makes
the unfolding precise and proves it well defined (Theorem~\ref{thm:unique}).

\begin{definition}[Semantics as solution]\label{def:semantics}
  The denotational semantics defined by a coalgebra \emph{is} its solution: state $x$ denotes
  $\behav(x)$.
\end{definition}

\noindent
This is a definition, not a theorem, and it is why no preservation result appears in this
paper: a solver that returns the solution returns the semantics, by definition. What remains
to prove is that the solution is well defined (Theorem~\ref{thm:unique}), that the solver
returns it (Theorem~\ref{thm:correctness}), and when the solver halts
(Theorems~\ref{thm:termination}, \ref{thm:optimal}).

\subsection{The Approximation Order}\label{sec:order}

Solutions are possibly infinite trees, so we fix the space they live in and the order along
which they are approximated.

\begin{definition}[Trees]\label{def:trees}
  $D$ is the set of finite and infinite trees in which every node is either an operation of
  $F$ with one subtree per argument, or a $\bot$ leaf, read ``nothing exposed here, yet or
  ever''. (Under a non-deterministic or probabilistic $M$, each node carries the
  $M$-structure of its layer; the deterministic reading suffices for this paper.)
\end{definition}

\begin{definition}[Approximation order]\label{def:order}
  For $t, u \in D$, $t \sqsubseteq u$ iff $u$ is obtained from $t$ by replacing $\bot$ leaves
  with subtrees. The single $\bot$ leaf is the least tree.
\end{definition}

\begin{proposition}[The tree space is a domain]\label{prop:order}
  $(D, \sqsubseteq)$ is a pointed $\omega$-complete partial order: every increasing chain has
  a least upper bound, computed level by level. Proof in Appendix~\ref{app:order}.
\end{proposition}

The order exists for one reason. $\bot$ is the value of ``nothing exposed'', the starting
point every iteration ascends from; reusing any exposed value as the starting point would
make a state that has not yet been computed indistinguishable from one that computed that
value.

\begin{definition}[Unfolding functional]\label{def:unfolding}
  The \emph{unfolding functional} $\Psi$ takes a candidate solution $g \colon X \to D$ to the
  candidate that reads one more layer off the coalgebra:
  \[
    \Psi(g)(x) =
    \begin{cases}
      \bot & \outm(x) = \bot,\\
      \sigma\big(g(x_1), \dots, g(x_k)\big) & \outm(x) \text{ exposes } \sigma(x_1, \dots, x_k).
    \end{cases}
  \]
  A solution of the coalgebra is exactly a fixpoint of $\Psi$.
\end{definition}

\begin{theorem}[The solution is unique]\label{thm:unique}
  $\Psi$ has exactly one fixpoint $\behav$, and it is simultaneously the order-theoretic least
  fixpoint $\bigsqcup_k \Psi^k(\lambda x.\bot)$ and the greatest, coinductive one. Proof in
  Appendix~\ref{app:unique}.
\end{theorem}

\noindent
The reason is guardedness: $\Psi$ commits to the root of every tree before consulting its
argument, so two candidate solutions can disagree only ever deeper, and under the metric in
which the distance between two trees halves with the depth of their first disagreement, $\Psi$
is a contraction, forcing its fixpoints equal. Consequently induction and coinduction coincide here, and \emph{any}
procedure that reaches a fixpoint reaches the solution, which is what licenses an interpreter
and a compiler with different evaluation orders.

\begin{remark}[Algebraic compactness]\label{rem:compact}
  Over discrete sets the least and greatest fixpoints of a signature differ (finite lists
  versus finite-or-infinite lists). Completing with $\bot$ and ordering by approximation is
  exactly what collapses them; Theorem~\ref{thm:unique} is the operational face of algebraic
  compactness \cite{smyth1982-recursive-domain-equations, freyd1991-algebraically-complete-categories}.
\end{remark}

\begin{remark}[Where a least fixpoint is essential]\label{rem:leastessential}
  Uniqueness holds for the unfolding, after $\outm$ has committed a root. Producing $\outm$
  itself from a reduction function is unguarded, and there least genuinely differs from
  greatest; that single least fixpoint lives in Section~\ref{sec:reductions}.
\end{remark}

\subsection{Tabling}\label{sec:tabling}

The solver computes $\behav$ at a root state by tabling $\outm$.

\begin{definition}[The solver]\label{def:solver}
  States are \emph{interned}: equal states are one object, so identity is one pointer test.
  The solver keeps a table from states to graph nodes. At a state $x$ it computes $\outm(x)$;
  if $\outm(x) = \bot$ the node is a $\bot$ leaf, and otherwise the node carries the exposed
  operation with one edge per successor, where a successor already tabled \emph{reuses} its
  node, and a reuse of a state still being processed is a \emph{back edge}, closing a cycle.
  The output is the table: a graph with one node per distinct reachable state.
\end{definition}

The solver consults only $\outm$ and the identity of states. It does not ask where the
coalgebra came from, and nothing in it is specific to any language.

\begin{theorem}[Soundness]\label{thm:correctness}
  The graph the solver builds at root $x_0$ unfolds to $\behav(x_0)$: the table is the
  restriction of the coalgebra to the states reachable from $x_0$, and its unfolding and
  $\behav$ are both fixpoints of the unfolding functional, hence equal by
  Theorem~\ref{thm:unique}. Proof in Appendix~\ref{app:correctness}.
\end{theorem}

\noindent
The two ingredients are worth naming. Reuse never changes the solution because the solver
reuses a node only for the \emph{identical} state, and identical states have identical
$\outm$, hence identical subtrees. And a back edge never changes the solution because the
finite cyclic graph and the infinite unfolding are the same fixpoint read two ways.

\begin{definition}[Rational]\label{def:rational}
  A tree is \emph{rational} if it has finitely many distinct subtrees; equivalently, it is the
  unfolding of a finite graph \cite{courcelle1983-infinite-trees}.
\end{definition}

\begin{theorem}[Termination]\label{thm:termination}
  When the silent run from each reachable state terminates, exposing a layer or returning to a
  visited state, so the solver computes $\outm$ there, the solver halts at $x_0$ iff finitely many
  states are reachable from $x_0$, and then its output is a finite cyclic graph unfolding to the
  rational solution $\behav(x_0)$. Proof in Appendix~\ref{app:termination}.
\end{theorem}

\begin{theorem}[Optimality]\label{thm:optimal}
  Fix the reduction function, and with it the solution it defines; a procedure is
  \emph{solution-preserving} when it returns that same solution as a finite graph. Such a graph
  exists iff the solution is rational, so no solution-preserving procedure halts on an input whose
  solution is not rational: the rational fragment is the ceiling. The bound is at this fixed
  solution: a procedure halting on a non-rational input while returning a finite graph returns a
  rational tree, hence not the solution fixed here, but a coarser semantics. The solver halts on
  the inputs with finitely many reachable states
  (Theorem~\ref{thm:termination}), a sub-fragment of the rational ones under structural state
  identity; the residual gap, rational solutions whose states are nonetheless pairwise distinct,
  is closed by folding behaviourally equal states, the coarsest solution-preserving identity,
  which attains the whole rational fragment but is in general undecidable. Proof in
  Appendix~\ref{app:optimal}.
\end{theorem}

\noindent
Rationality is undecidable, so the halting inputs cannot be pre-selected. The solver is a sound
semi-decision procedure: it returns the finite graph on the inputs with finitely many reachable
states and diverges on the rest, and outside the rational fragment no solution-preserving
procedure halts at all, since the solution it would return as a finite graph does not exist.

\subsection{Reduction Functions Present Coalgebras}\label{sec:reductions}

Everything so far assumed $\outm$ given. The coalgebras of interest are presented differently:
by a \emph{rule} that at each state either exposes a layer or takes an internal step, and only
finitely many internal steps separate a state from its next layer, if one comes at all.

\begin{definition}[Reduction function]\label{def:reduction}
  A \emph{reduction function} is a map
  \[
    \stepm \colon X \longrightarrow M\big(F X + X\big):
  \]
  at each state, one $M$-step that either exposes a layer of data (left summand) or continues
  silently at a successor state (right summand). For this definition $M$ must support
  iteration: each $M A$ is a pointed $\omega$-complete partial order with continuous
  composition, as $\Delaym$ and the other effects of Definition~\ref{def:signature} are
  \cite{adamek2006-iterative-algebras, goncharov2016-complete-elgot-monads}.
\end{definition}

\begin{definition}[The structure map]\label{def:out}
  The reduction function $\stepm$ \emph{presents} the coalgebra $\outm = \iter{\stepm}$, the
  least fixpoint of
  \[
    \Phi(g)(x) = \big[\,\text{expose},\; g\,\big]\big(\stepm(x)\big),
  \]
  where the case split passes an exposed layer through unchanged and applies $g$ to a silent
  successor. The fixpoint runs silent steps until the first exposed layer: $\outm(x)$ is that
  layer, and where the silent process from $x$ never exposes one, $\outm(x) = \bot$.
\end{definition}

\begin{proposition}[The presentation is well defined]\label{prop:out}
  $\Phi$ is monotone and $\omega$-continuous, so its least fixpoint
  $\outm = \bigsqcup_k \Phi^k(\lambda x.\bot)$ exists, and $\outm(x) = \bot$ exactly when the
  silent process from $x$ exposes no layer. Proof in Appendix~\ref{app:out}.
\end{proposition}

\noindent
This is the one least fixpoint of the paper, and the one place least differs from greatest: a
silent loop exposes nothing, so its value is forced to $\bot$. The solver computes $\outm(x)$
by running the silent steps, and the states a silent run passes through are interned like all
others, so the run keeps the set of states it has visited and stops at a revisit. Because the
deterministic silent step is a function, a revisited state replays the same run forever and
exposes no layer, so returning $\bot$ at the revisit is exactly the least fixpoint, decided in
finite time. Everything above now applies to
$(X, \outm)$ unchanged. Conversely no generality was lost by developing the theory for
coalgebras: a coalgebra is the degenerate reduction function with no silent steps, so the two
formats present the same objects.

Two independent finiteness conditions now govern the solver: the silent run from each visited
state must reach a layer or a revisit (this section), and the reachable states must be finite
(Theorem~\ref{thm:termination}); either can fail without the other.

For instance, one step of weak-head reduction in the pure $\lambda$-calculus is such a
reduction function: a variable or an abstraction exposes a layer, and an application either
contracts its head redex, the abstraction at the head of its function spine applied to the next
argument, as a silent step, or exposes a layer when that head is a variable. Its solution is the L\'evy--Longo tree~\cite{longo1983-set-theoretical-models} of the
term, the standard tree semantics of weak-head reduction; a different reduction function fixes a
different such tree (the B\"ohm tree under head reduction), and under
Definition~\ref{def:semantics} each \emph{is} the meaning its reduction defines. Here a bare
coalgebra has only a solution, no prior meaning to preserve, but a programming language carries a
standard denotation of its terms; when that denotation is the tree the chosen reduction unfolds
to, as the L\'evy--Longo tree is for weak-head reduction, the solution-preserving guarantee of
Theorem~\ref{thm:optimal} is exactly that this denotational semantics is unchanged. The instance
is developed in Section~\ref{sec:bridge}.

% =====================================================================
\section{Implementation}\label{sec:bridge}
% =====================================================================
The method is independent of the $\lambda$-calculus, but its first instance is the
$\lambda$-calculus itself, and we have built it: an interpreter for weak-head reduction and, over
it, a self-hosting compiler. Both accompany the paper\anon[ as supplementary
material]{~\cite{colambda}}, and every behaviour stated below is reproduced as a passing test.

\paragraph{The bridge.}
The signature is $F X = \{\Varsh_i\}_i + \{\Lamsh(x)\}_x + \{\Appsh(x, x')\}_{x, x'}$, a nullary
operation per de Bruijn index, a unary $\Lamsh$, a binary $\Appsh$; the states are pure
$\lambda$-terms in de Bruijn form, with no recursion binder. They are \emph{interned}, so two
structurally equal terms are one object and identity is a pointer test, the decidable, hashable
state identity Definition~\ref{def:coalgebra} asks for and the reason the interpreter decides on
the syntax of a term what $\beta$-equality, the identity of the function it computes, cannot. The
reduction function is one step of weak-head reduction: a variable or an abstraction exposes its
layer, and an application either contracts its head redex, the abstraction at the head of its
function spine applied to the next argument, as a silent step, or exposes an $\Appsh$ layer when
that head is a variable; the silent iteration $\outm = \iter{\stepm}$ then runs these contractions
to the first layer, or to $\bot$ when the head never settles.

\begin{proposition}[The weak-head instance]\label{prop:bridge}
  For this $\stepm$, the structure map $\outm = \iter{\stepm}$ is the weak-head shape, with
  $\outm(x) = \bot$ exactly when $x$ has no weak head normal form, so its solution $\behav$ is by
  definition the term's L\'evy--Longo tree~\cite{longo1983-set-theoretical-models}, the standard
  lazy tree semantics. A second structure map that reduces under $\lambda$, reading a $\lambda$
  over a body whose own solution is $\bot$ as $\bot$, gives the B\"ohm tree instead. Both
  specialise Propositions~\ref{prop:order}, \ref{prop:out} and Theorem~\ref{thm:correctness} to
  $M = \Delaym$, so the general solver computes each.
\end{proposition}

\paragraph{The engine.}
The interpreter realises the two least fixpoints of Section~\ref{sec:method} as two mechanisms on
the interned term. The silent collapse $\outm$ is a memoised recursion that runs head reduction
and, on finding a state already on its own stack by pointer identity, returns $\bot$; so $\Omega$
is decided as $\bot$ in finite time where reduction loops. The readout is the interning itself:
successor states are the de Bruijn reducts the reduction produces, so when a reduct is structurally
identical to a live state the back edge is there already and the cycle folds at no cost. Thus
$Y\,(\mathtt{cons}\,0)$, whose self-application reduces to a structurally identical term, comes out
as a finite cyclic graph, a cell whose tail points back to itself, and the ordinary
$\mathtt{map}\,\mathtt{succ}$ over it as the finite circle of ones, with nothing cycle-aware in
either term; a term whose reducts never coincide, such as the $Y\,F\,\underline{0}$ of
Appendix~\ref{app:optimal}, does not fold, exactly the structural-identity gap proved there. The
weak-head and B\"ohm readings are two structure maps over the one carrier, the head reading walking
the same states the weak-head readout does, so the solver halts on the same terms under both,
differing only at a $\lambda$ over a body with no head normal form.

\paragraph{Compilation as a second coalgebra.}
Compilation is itself such an equation: it sends each sub-term of a quoted source program, by
specialisation rules, to a node of a Python abstract syntax tree over that sub-term's own
sub-terms, and its solution is the program compiled to that tree. The compiler is a pure
$\lambda$-term the interpreter solves on the quoted source, and so are its enabling analyses, a
simple-typability check by Hindley--Milner inference with an occurs check, a closedness check, and
a fuel-bounded normalisation check, so program analysis is expressed as evaluation. The analyses
select, per sub-term, one of three targets: a closed simply-typable sub-term, hence strongly
normalising, compiles to a strict call-by-value target; one whose interpreted behaviour is a
finite normal form, but untypable, to a call-by-name target; everything else stays interpreted,
the interpreter being the sound default. A compiled sub-term is a single node holding native code,
spliced into the otherwise interpreted graph, so specialisation is local. The fixpoint combinator
the compiler emits is a strict, call-by-value variant $Z$ rather than $Y$, so the compiled Python
runs the recursion under call-by-value, where a compiled $Y$ would diverge at once, and the
compiler compiles its own source and self-hosts; the cost of the multi-stage self-compilation is the subject of
Section~\ref{sec:application}.

% =====================================================================
%  TODO: REMAINING SECTIONS (paper outline, IMRaD-style; sections 1-2 done)
% ---------------------------------------------------------------------
%  Overall structure (agreed with author). This paper is not pure theory:
%  it ships an interpreter, a self-hosting compiler, and reproducible
%  experiments, so an IMRaD shape fits its actual content better than the
%  legacy "theory-to-the-end, engineering-in-appendices" layout.
%
%    1. Introduction        -- DONE (difficulty stack -> resolution -> the
%                              coalgebra/reduction/lambda syllogism).
%    2. The Method          -- DONE (coalgebras; approximation order;
%                              tabling; reduction functions present them).
%                              Proofs in appendices app:order/out/unique/
%                              correctness/termination/optimal.
%    3. Implementation      -- TODO. OPENS with the lambda bridge (the
%                              \ref{sec:bridge} instance), then the engine.
%                              (a) The bridge: weak-head reduction is one e;
%                              its solution is the Levy-Longo tree (= the
%                              standard meaning); lambda signature
%                              Var/Lam/App, interning on de Bruijn terms.
%                              Theorem-grade content (cf. legacy
%                              prop:shisout), and the rung from the abstract
%                              method to the concrete engine. (b) The
%                              interpreter (interning; the two
%                              least-fixpoint mechanisms). (c) The compiler /
%                              self-hosting bootstrap (eager / call-by-name /
%                              fixpoint targets; closed-subterm island
%                              specialization; the strict Z combinator
%                              bootstrap). Source: legacy app:compiler,
%                              moved up out of the appendix per the author's
%                              "the compiler does not belong in an appendix".
%    4. Evaluation          -- TODO. (NOT "Results": PL venues expect
%                              "Evaluation" / "Case Studies".) Measured
%                              behaviour: Y(cons 0) folds to a finite cyclic
%                              graph; Omega and Y(\x.x) decided as bot; ed
%                              runs in O(mn); the DSL case studies
%                              (Datalog / reachability+points-to / minimax /
%                              lexer, from app:dsl, results cited, lambda
%                              encodings left in the appendix); the
%                              self-compilation benchmark matrix. State the
%                              flat-matrix fact here plainly.
%    5. Discussion          -- TODO. Where legacy's scattered remarks find a
%                              home as a coherent limitations/analysis
%                              section: the two independent finiteness
%                              conditions; optimal semi-decision and the
%                              undecidability of rationality; optimality is
%                              relative to a fixed denotation; the
%                              structural-identity vs behavioural-equivalence
%                              gap; WHY the benchmark matrix is flat (an
%                              island gives up the shared interner cache:
%                              local soundness vs global performance);
%                              comparison with call-by-need.
%    6. Related work        -- TODO. Rewrite of legacy sec:related. Deferred
%                              to here (not before the method) so each
%                              comparison (CoPaR / partition refinement,
%                              tabling & least Herbrand models, graph
%                              reduction & cyclic lambda, rational trees &
%                              recursive types, higher-order model checking)
%                              lands after the construction is in hand.
%    7. Conclusion          -- TODO. Rewrite.
%
%  The lambda bridge opens section 3 (author's decision): the Levy-Longo
%  theorem is stated there, before the engine that realises it, so the jump
%  from the abstract method (section 2, lambda-free) to the concrete
%  implementation passes through the instance. Section 2 therefore ends at
%  full generality; \ref{sec:bridge} now resolves into section 3.
%
%  HARD CONSTRAINTS carried from sections 1-2: every legacy line is opt-in
%  (rewrite, never copy); the general layers (method, evaluation of the
%  generality claim) speak of any coalgebra / any reduction function, with
%  the lambda-calculus appearing only as an explicit instance; no dashes in
%  prose; no conversational meta-vocabulary (flagship/reveal/hook) in the
%  prose; blind-read each new section to a fresh context-free reader.
% =====================================================================

\bibliographystyle{ACM-Reference-Format}
\bibliography{references}

\appendix

% =====================================================================
\section{The Tree Space Is a Domain}\label{app:order}
% =====================================================================
\begin{proof}[Proof of Proposition~\ref{prop:order}]
  Present a tree $t \in D$ as a partial labelling: a prefix-closed set of positions, each
  labelled by an operation of $F$ or by $\bot$, a node labelled by an operation of arity $k$
  having exactly $k$ children and a $\bot$ node having none.

  \emph{Least element.} The single $\bot$ leaf has the one position $\varepsilon$, labelled
  $\bot$, and every tree refines it, so it is least under Definition~\ref{def:order}.

  \emph{Least upper bounds.} Let $t_0 \sqsubseteq t_1 \sqsubseteq \cdots$ be a chain. Define
  $t$ on the union of their positions, labelling a position with the operation any $t_i$
  assigns it, and with $\bot$ if every $t_i$ that contains the position assigns $\bot$. The
  operation label is well defined: refinement never relabels an operation, only replaces
  $\bot$, so two members of the chain agree on every position they both label with
  operations. The position set of $t$ is prefix-closed as a union of prefix-closed sets, and
  arity-respecting because each label comes with its children from some $t_i$. Each
  $t_i \sqsubseteq t$ by construction, and any upper bound $u$ refines every $t_i$, hence
  carries every operation label of $t$, so $t \sqsubseteq u$. Thus $t$ is the least upper
  bound, and it is computed level by level: the label of a position in $t$ is determined by
  any member of the chain that already exposes it.
\end{proof}

% =====================================================================
\section{The Presentation Is Well Defined}\label{app:out}
% =====================================================================
\begin{proof}[Proof of Proposition~\ref{prop:out}]
  \emph{Monotone and continuous.} $\Phi(g)(x)$ post-composes the fixed computation
  $\stepm(x)$ with the case split $[\text{expose}, g]$, which depends on $g$ only in the
  silent branch. Composition in $M$ is continuous and the case split preserves it, so $\Phi$
  is monotone and $\omega$-continuous in the pointwise order on $X \to M(F X)$, whose
  structure is Proposition~\ref{prop:order}'s argument applied to $M(F X)$. By Kleene,
  $\outm = \bigsqcup_k \Phi^k(\lambda x.\bot)$ exists and is the least fixpoint of $\Phi$.

  \emph{Characterisation.} We show, for the deterministic effect $M = \Delaym$ and by
  induction on $k$: $\Phi^k(\lambda x.\bot)(x) \neq \bot$ iff the silent process from $x$
  exposes a layer within fewer than $k$ steps, and then $\Phi^k(\lambda x.\bot)(x)$ is that
  layer.

  For $k = 0$ both sides are false: $\Phi^0(\lambda x.\bot)(x) = \bot$, and no layer is
  exposed within fewer than zero steps.

  For $k + 1$, unfold one step: $\Phi^{k+1}(\lambda x.\bot)(x) =
  [\text{expose}, \Phi^k(\lambda x.\bot)](\stepm(x))$. If $\stepm(x)$ exposes a layer, the
  value is that layer, which is non-$\bot$, and the layer is exposed within zero steps. If
  $\stepm(x)$ continues silently at $x'$, the value is $\Phi^k(\lambda x.\bot)(x')$, which by
  the induction hypothesis is non-$\bot$ iff $x'$ exposes within fewer than $k$ steps, that
  is, iff $x$ exposes within fewer than $k + 1$ steps, and is then that layer.

  Taking the least upper bound over $k$: $\outm(x) \neq \bot$ iff the silent process from $x$
  exposes a layer within some finite number of steps, and $\outm(x)$ is then that layer;
  $\outm(x) = \bot$ exactly when no layer is ever exposed. This characterisation is the
  deterministic case the paper uses; for a richer effect the iteration is the standard one of
  complete Elgot monads \cite{adamek2006-iterative-algebras, goncharov2016-complete-elgot-monads},
  which we do not develop here.
\end{proof}

% =====================================================================
\section{Uniqueness of the Solution}\label{app:unique}
% =====================================================================
\begin{definition}[Tree metric]\label{def:metric}
  For $t, u \in D$ let $\mathrm{d}(t, u) = 0$ if $t = u$ and otherwise
  $\mathrm{d}(t, u) = 2^{-\ell}$, where $\ell$ is the least depth at which $t$ and $u$
  differ. Lift it to candidate solutions by
  $\mathrm{d}(g, h) = \sup_x \mathrm{d}(g(x), h(x))$.
\end{definition}

\begin{lemma}[Completeness]\label{lem:metric}
  $(D, \mathrm{d})$ and the candidate space $(X \to D, \mathrm{d})$ are complete metric
  spaces, and the limit of an increasing chain that converges in $\mathrm{d}$ is its least
  upper bound.
\end{lemma}
\begin{proof}
  A Cauchy sequence of trees is eventually constant at every fixed depth, so the levelwise
  limit exists in $D$ and the sequence converges to it. The candidate space inherits this:
  a sequence Cauchy in the supremum metric is uniformly Cauchy in $x$, so it is pointwise
  Cauchy with a modulus independent of $x$; its pointwise limit $g$ therefore satisfies
  $\sup_x \mathrm{d}(g_n(x), g(x)) \to 0$ (let $m \to \infty$ in the uniform Cauchy bound),
  so $g_n \to g$ in $\mathrm{d}$ and the space is complete. For an
  increasing chain, the levelwise limit and the levelwise least upper bound of
  Proposition~\ref{prop:order} are the same tree.
\end{proof}

\begin{lemma}[Guardedness]\label{lem:guarded}
  $\mathrm{d}(\Psi(g), \Psi(h)) \le \tfrac{1}{2}\,\mathrm{d}(g, h)$.
\end{lemma}
\begin{proof}
  Fix $x$. The root of $\Psi(g)(x)$ is determined by $\outm(x)$ alone: $\bot$, or the exposed
  operation. So $\Psi(g)(x)$ and $\Psi(h)(x)$ share their root, and they can differ only
  inside some child, where the difference is one of $g(x_i)$ against $h(x_i)$, one level
  deeper. Hence the least depth of difference grows by one, halving the distance.
\end{proof}

\begin{proof}[Proof of Theorem~\ref{thm:unique}]
  By Lemma~\ref{lem:guarded}, $\Psi$ is a contraction on the complete space of
  Lemma~\ref{lem:metric}, so by Banach's fixed-point theorem \cite{banach1922-operations} it has exactly
  one fixpoint $\behav$, and every iteration sequence converges to it. The particular
  sequence $\Psi^k(\lambda x.\bot)$ is increasing, since $\lambda x.\bot$ is least and $\Psi$
  is monotone, and converges to $\behav$ in $\mathrm{d}$; by Lemma~\ref{lem:metric} its least
  upper bound is that limit, so the order-theoretic least fixpoint equals $\behav$. The
  coinductive unfolding, by its defining property of reading layers off $\outm$, is a
  fixpoint of $\Psi$, hence also $\behav$.
\end{proof}

% =====================================================================
\section{Soundness of the Solver}\label{app:correctness}
% =====================================================================
\begin{lemma}[The table is the reachable restriction]\label{lem:table}
  Let $R \subseteq X$ be the least set containing the root $x_0$ and closed under successors
  of $\outm$. The solver's table assigns exactly one node to each state of $R$, the node
  carrying $\outm(x)$: interning makes equal states one table key, the solver computes
  $\outm(x)$ once per key, and each edge of the node goes to the node of a successor state.
  In particular $(R, \outm)$ is itself a coalgebra: successors of reachable states are
  reachable.
\end{lemma}

\begin{proof}[Proof of Theorem~\ref{thm:correctness}]
  Let $u \colon R \to D$ assign to each tabled state the unfolding of its graph node. By
  Lemma~\ref{lem:table}, $u(x)$ has the root $\outm(x)$ dictates and the subtrees
  $u(x_1), \dots, u(x_k)$ of the successor nodes; this is exactly the statement that $u$ is a
  fixpoint of the unfolding functional of the coalgebra $(R, \outm)$. A back edge changes
  nothing in this reading: the reused node is the node of the identical state, whose
  unfolding is the same subtree. The restriction $\behav|_R$ is also a fixpoint of that
  functional, because $\behav$ reads layers off the same $\outm$ and $R$ is closed under
  successors. By Theorem~\ref{thm:unique}, applied to the coalgebra $(R, \outm)$, the two
  fixpoints are equal, so $u(x_0) = \behav(x_0)$: the graph unfolds to the solution.
\end{proof}

% =====================================================================
\section{Termination of the Solver}\label{app:termination}
% =====================================================================
\begin{proof}[Proof of Theorem~\ref{thm:termination}]
  Suppose $R$, the states reachable from $x_0$, is finite. Interning gives each state one
  table key, so the solver computes $\outm$ at most once per state of $R$, each call halting
  because the silent run from its state terminates by hypothesis, and follows finitely many
  edges per state; it therefore halts. Its output
  unfolds to $\behav(x_0)$ by Theorem~\ref{thm:correctness}, and the unfolding has at most
  one distinct subtree per node of the finite graph, so it is rational
  (Definition~\ref{def:rational}).

  Conversely, the solver tables every reachable state before it can return, one node per
  state by Lemma~\ref{lem:table}; if $R$ is infinite, the table never closes and the solver
  diverges.
\end{proof}

% =====================================================================
\section{Optimality of the Solver}\label{app:optimal}
% =====================================================================
\begin{proof}[Proof of Theorem~\ref{thm:optimal}]
  \emph{A solution has a finite representation iff it is rational.} The unfolding of a finite
  graph has at most one distinct subtree per node, hence finitely many, so it is rational.
  Conversely a rational tree is the unfolding of the finite graph whose nodes are its
  distinct subtrees, each node carrying its root operation with an edge to each immediate
  subtree; this is the standard correspondence for regular trees \cite{courcelle1983-infinite-trees}.

  \emph{Optimality.} Fix the reduction function and the solution it defines. A
  \emph{solution-preserving} procedure returns this solution as a finite graph; a finite
  generator such as the input term names the solution but does not return it as a graph, and is
  not of this kind. By the first part such a graph exists only when the solution is rational, so a
  solution-preserving procedure cannot halt on an input whose solution is not rational: the
  rational fragment bounds them all. The bound is at the fixed solution: to halt on a non-rational
  input while returning a finite graph one must return a rational tree, hence not the non-rational
  solution fixed here; such a procedure computes a coarser semantics, a different reduction's
  solution, not this one. The solver halts exactly on inputs with finitely many reachable states
  (Theorem~\ref{thm:termination}), all of whose solutions are rational, and its output is that
  finite graph (Theorem~\ref{thm:correctness}). Within the rational fragment a gap remains under
  structural state identity, witnessed in the $\lambda$-instance by $Y\,F\,\underline{0}$ with
  $F = \lambda s.\lambda n.\,\mathtt{cons}\,0\,(s\,(\mathtt{succ}\,n))$, whose states
  $Y\,F\,\underline{k}$ for $k = 0, 1, 2, \dots$ are pairwise distinct yet all unfold to the one
  rational stream of zeros, so structural folding never closes though the solution is rational and
  the solver diverges; the solver, folding on structural identity, halts on the rational inputs
  outside that gap. The coarsest solution-preserving identity folds states with equal subtrees,
  whose classes are the distinct subtrees of the solution, finite exactly when the solution is
  rational; that identity therefore attains the whole rational fragment, but it is in general
  undecidable (on $\lambda$-terms it is behavioural equality of states), so the solver's decidable
  structural identity attains only a sub-fragment. We do not claim the rational fragment is
  unattainable by every decidable identity, only that the structural one the solver uses falls
  short of it.
\end{proof}

\section{The Rational Fragment as a Domain-Specific Language}\label{app:dsl}
The terminating fragment of the construction, the rational behaviours, is exactly the world of
finite-state, regular, and least-fixpoint-over-a-finite-lattice computation, and a range of
practical domains live there. Each domain below is a pure $\lambda$-term that reaches finitely
many states, so its behaviour is rational and the interpreter halts, and each is reproduced as a
test\anon[~(supplementary
material)]{~\cite{colambda}}. Two reusable shapes recur: a least fixpoint over a Church
tuple of booleans (the Datalog engine of Section~\ref{sec:application}), and a recursion whose
subproblems are the nodes of a first-order datum, memoised for free because identical subproblems
are one interned state. Each claim is confined to a domain's rational core: regular lexing but not
general context-free parsing, rational-tree unification but not a general most-general unifier,
finite-state model checking but not infinite-state. The terms below are the actual encodings; the
Booleans are Church, $\mathtt{T} = \lambda a.\lambda b.\,a$ and $\mathtt{F} = \lambda a.\lambda b.\,b$,
with $\mathtt{or} = \lambda p.\lambda q.\,p\,\mathtt{T}\,q$ and
$\mathtt{and} = \lambda p.\lambda q.\,p\,q\,\mathtt{F}$, and $\mathtt{cons}$ and $Y$ are as in
Section~\ref{sec:bridge}.

\paragraph{Mapping a cyclic list.}
The $\mathtt{map}$ of the introduction is the ordinary one, with nothing cycle-aware in it,
\[
  \mathtt{map}\,f = Y\,\big(\lambda\mathit{self}.\,\lambda\mathit{lst}.\;
  \mathit{lst}\,(\lambda h.\lambda t.\,\mathtt{cons}\,(f\,h)\,(\mathit{self}\,t))\;\mathtt{nil}\big),
\]
where a Scott-encoded list applies one of two handlers, for a $\mathtt{cons}\,h\,t$ cell or for
the empty list $\mathtt{nil}$. Applied to $r = Y\,(\mathtt{cons}\,0)$, the
recursive call $\mathit{self}\,t$ re-enters the same interned state, so tabling folds it and the
result is a finite circular list, a single cell holding $f\,0$ whose tail points back to itself,
every element transformed, where ordinary reduction would map the infinite stream forever. On a finite list the
same $\mathtt{map}$ is the textbook one.

\paragraph{Dynamic programming, memoisation for free.}
The decisive capability is memoisation: a recursion whose state space is a
tree or graph computes each distinct subproblem once, because a structurally repeated subtree, or a
graph node revisited along several paths, is one interned state. A perfect binary tree of depth $n$
whose nodes share both children is a graph of $n+1$ states that unfolds to $2^{n}$ leaves; a tree
DP over it is read in time linear in $n$, where a naive recursion would visit $2^{n}$ leaves. No
memo table is written; tabling on first-order state identity is the memoisation, which the pure
$\lambda$-calculus lacks without a decidable identity on subproblems, value equality of closures
being undecidable.
\[
  \mathtt{node} = \lambda l.\lambda r.\lambda m.\lambda f.\,m\,l\,r, \qquad
  \mathtt{leaf} = \lambda v.\lambda m.\lambda f.\,f\,v,
\]
\[
  \mathtt{any} = Y\,\big(\lambda\mathit{self}.\lambda t.\;
  t\,(\lambda l.\lambda r.\,\mathtt{or}\,(\mathit{self}\,l)\,(\mathit{self}\,r))\,(\lambda v.\,v)\big),
  \qquad t_0 = \mathtt{leaf}\;\mathtt{F},\quad t_{k+1} = \mathtt{node}\;t_k\;t_k,
\]
so $\mathtt{any}\,t_n$ reads to $\mathtt{F}$ with each of the $n+1$ shared subtrees evaluated once.

\paragraph{Game search with a transposition table.}
Minimax is an and-or tree: a maximising position takes the disjunction over its moves, a minimising
position the conjunction, and a terminal position a Boolean. A \emph{transposition}, the same
position reached by a different move order, is the same interned node, so its value is computed
once. Interning is the transposition table of a game-playing program, with no table written by hand.
\[
  \mathtt{max} = \lambda l.\lambda r.\lambda x.\lambda n.\lambda f.\,x\,l\,r, \quad
  \mathtt{min} = \lambda l.\lambda r.\lambda x.\lambda n.\lambda f.\,n\,l\,r, \quad
  \mathtt{leaf} = \lambda v.\lambda x.\lambda n.\lambda f.\,f\,v,
\]
\[
  \mathtt{value} = Y\,\big(\lambda\mathit{self}.\lambda p.\;
    p\,(\lambda l.\lambda r.\,\mathtt{or}\,(\mathit{self}\,l)(\mathit{self}\,r))\,
  (\lambda l.\lambda r.\,\mathtt{and}\,(\mathit{self}\,l)(\mathit{self}\,r))\,(\lambda v.\,v)\big).
\]

\paragraph{Graph reachability and program analysis.}
Reachability and transitive closure over a directed graph, including one with cycles, are the least
fixpoint of the immediate-consequence operator, terminating because the lattice is finite; the same
shape decides reachability of a state in a finite transition system. Andersen-style points-to (alias)
analysis \cite{andersen1994-program-analysis-c, bravenboer2009-doop}, in which $\mathtt{pointsTo}(p,o)$ follows from
$\mathtt{assign}(p,q)$ and $\mathtt{pointsTo}(q,o)$, is monotone Datalog, and a compiler reads it as
the same boolean fixpoint. The cyclic graph $a\to b\to c\to a$ with $c\to d$, and the points-to
program for $a=\mathtt{new}\,o_1;\,b=a;\,c=b$, are the clause sets
\[
  \begin{aligned}
    &\mathtt{reach}(a),\quad \mathtt{reach}(b)\leftarrow\mathtt{reach}(a),\quad
    \mathtt{reach}(c)\leftarrow\mathtt{reach}(b),\\
    &\mathtt{reach}(a)\leftarrow\mathtt{reach}(c),\quad \mathtt{reach}(d)\leftarrow\mathtt{reach}(c);\\
    &\mathtt{pt}(a,o_1),\quad \mathtt{pt}(b,X)\leftarrow\mathtt{pt}(a,X),\quad
    \mathtt{pt}(c,X)\leftarrow\mathtt{pt}(b,X),
  \end{aligned}
\]
so $\mathtt{reach}(d)$ and $\mathtt{pt}(c,o_1)$ hold while $\mathtt{reach}(e)$ and $\mathtt{pt}(c,o_2)$
do not.

\paragraph{Lexical analysis.}
A deterministic finite automaton is a finite cyclic transition graph, and running it is a fold over
the input that threads the state. Regular languages are the archetypal rational behaviour, so a lexer
is a $\lambda$-term whose behaviour is rational; the name \emph{rational} is itself the regular, or
Kleene, sense (Definition~\ref{def:rational}). Encoding the even state and the symbol $a$ as
$\mathtt{T}$, and the odd state and $b$ as $\mathtt{F}$, both two-way selectors,
\[
  \delta = \lambda s.\lambda x.\,s\,(x\,\mathtt{F}\,\mathtt{T})\,(x\,\mathtt{T}\,\mathtt{F}), \qquad
  \mathtt{accepts} = \lambda s.\,s\,\mathtt{T}\,\mathtt{F},
\]
\[
  \mathtt{run} = Y\,\big(\lambda\mathit{self}.\lambda s.\lambda w.\;
  w\,(\lambda h.\lambda t.\,\mathit{self}\,(\delta\,s\,h)\,t)\,s\big),
\]
and $\mathtt{accepts}\,(\mathtt{run}\,\mathtt{T}\,w)$ is $\mathtt{T}$ exactly when $w$ has an even
number of $a$s.

\paragraph{Unification and recursive types.}
When the solver halts on two terms, returning their finite graphs, equality of the behaviours
those graphs unfold to is decided by cycle-detecting bisimulation on the
graphs~\cite{courcelle1983-infinite-trees}. This is the decidable core shared by rational-tree
unification \cite{colmerauer1984-trees} and equirecursive ($\mu$-) type equality
\cite{amadio1993-subtyping-recursive-types, brandt1998-coinductive-type-equality}: two constructions of one rational term read out
identically, and a different term reads out differently. For instance $Y\,(\mathtt{cons}\;0)$ and
$Y\,(\lambda s.\,\mathtt{cons}\;0\;s)$ read out identically, while $Y\,(\mathtt{cons}\;1)$ differs.

\paragraph{The metalanguage needs tabling too.}
The interpreter's own host-language term builder illustrates the point reflexively. Writing a shared
subterm by reusing a builder gives a term whose construction unfolds the sharing exponentially, and
a cyclic builder would not terminate, until the builder is memoised by binder depth, the
host-language analogue of tabling. Intuitive code is slow or non-terminating exactly where the
metalanguage does not table. The memo key here is the syntactic builder rather than the first-order
behavioural identity the calculus folds on, but the moral is the same.

\paragraph{Further instances.}
The same rational core underlies others not reproduced here: equality saturation over e-graphs
\cite{willsey2021-egg, zhang2023-egglog} and binary decision diagrams \cite{bryant1986-bdd}, both of which are interning
of a shared graph; garbage-collection reachability over a cyclic heap; abstract interpretation as a
fixpoint over a finite abstract domain; and, with a non-trivial effect monad, the stationary
distribution of a finite Markov chain. Each is a finite-state, regular, or finite-lattice-fixpoint
computation, hence a rational behaviour the construction reads.

\end{document}
